\chapter{Δευτερο κεφάλαιο-2η Συνάντηση}
\label{ch:chapter2}


\section{\texorpdfstring{Διερεύνηση συσχέτισης $P-\cos\phi$ \textlatin{S} κυμματομορφών στις \textlatin{CDF}}{Διερεύνηση συσχέτισης P-cos(phi) S κυματομορφών στις CDF}}

Παρακάτω θα γίνει μια ανάλυση, ώστε να διερευνηθεί πιθανή συσχέτιση του $\cos\phi$ ως προς την ενεργή ισχύ $P$  σε κατανομές όπου παρατηρήθηκε ότι, ενώ τα ΦΒ παρήγαν ισχύ, δεν λειτουργούσαν με συντελεστή ισχύος μονάδα,αλλά το μεγαλύτερο μέρος της λειτουργίας τους στο εύρος $0.94$--$0.98$.

Tα παραπάνω ΦΒ παρουσιάζουν μια \textlatin{S} κυμματομορφή (σιγμοειδή) στη \textlatin{CDF}. Δηλαδή,η μεταβολή του $\cos\phi$ ως προς την ενεργή ισχύ $P$ ακολουθεί μη γραμμική συμπεριφορά.

\begin{figure}[H]
  \centering
  \includegraphics[width=0.82\textwidth]{δικτυο_paper_s_κυμματομορφων.png}
  \caption{Τα ΦΒ με  \textcolor{red}{\textlatin{S}} κυμματομορφή του  δίκτυο διανομής υπό εξέταση (πηγή:~\cite{dimoulias2024droop}).}
  \label{fig:network_s_waveform_paper}
\end{figure}

\subsection{Επιλογή ηλιόλουστης ημέρας μεγιστης παραγωγής ισχύος }


Αρχικά, η ανάλυση συσχέτισης πραγματοποιείται σε μία αντιπροσωπευτική ηλιόλουστη ημέρα, για την οποία η καμπύλη παραγωγής των ΦΒ εμφανίζει μονοκορυφική μορφή ή αλλιώς <<καμπάνα>>. Ο ορισμός της ημέρας αυτής βασίζεται στα ακόλουθα κριτήρια:

\begin{enumerate}
  \item \textbf{Κριτήριο μέγιστης ημερήσιας ενέργειας.} Για κάθε ημέρα $d$ υπολογίζεται η ημερήσια ενέργεια $E_d$ (σε \textlatin{kWh}) από τη χρονοσειρά ενεργής ισχύος $P_d(t_k)$ (σε \textlatin{kW}):
  \begin{equation}
    E_d \approx \sum_{k=1}^{N_d} P_d(t_k)\,\Delta t,
  \end{equation}
  όπου $\Delta t$ το χρονικό βήμα δειγματοληψίας ($15 \textlatin{min}$ ή $15 $ λεπτά - $0.25\textlatin{h}$).\\ Επομένως, η $E_d$ υπολογίζεται σε \textlatin{kWh}. Επιλέγεται η ημέρα
  \begin{equation}
    d^{\star}=\arg\max_d E_d.
  \end{equation}
Δηλαδή η ημέρα με την μέγιστη παργαγωγή ενεργής ισχύος  \item \textbf{Κριτήριο μορφής καμπύλης ισχύος.} Για την επιλεγμένη ημέρα $d^{\star}$ απαιτείται η καμπύλη ισχύος να έχει ένα ακρότατο και αυτό να είναι το ολικό μέγιστο, δηλαδή να υπάρχει χρονική στιγμή κορυφής $t_{\mathrm{peak}}$ τέτοια ώστε:
  \begin{equation}
    t_i<t_j\le t_{\mathrm{peak}} \Rightarrow P(t_i)<P(t_j)
    \quad (\text{γνησίως αυξούσα}),
  \end{equation}
  \begin{equation}
    t_{\mathrm{peak}}\le t_i<t_j \Rightarrow P(t_i)>P(t_j)
    \quad (\text{γνησίως φθίνουσα}).
  \end{equation}
\end{enumerate}

Συνοπτικά, η επιθυμητή συμπεριφορά είναι:
\begin{equation}
  P(t): \text{πρωί} \nearrow P_{\max} \searrow \text{απόγευμα},
\end{equation}
ώστε η παραγωγή να προσεγγίζει τη χαρακτηριστική «καμπάνα» καθαρών συνθηκών ακτινοβολίας.

Με βάση τα παραπάνω κριτήρια, επιλέχθηκε μία αντιπροσωπευτική ηλιόλουστη ημέρα του Μαΐου  και συγκεκριμένα η ημερομηνία \textbf{(2024-05-06)}.
 \begin{figure}[H]
  \centering
  \includegraphics[width=0.82\textwidth]{PICS/116310-11.png}
  \caption{Ενδεικτική ημερήσια καμπύλη ενεργού ισχύος για το ΦΒ 116310-\textbf{11} κατά την επιλεγμένη ηλιόλουστη ημέρα.}
  \label{fig:pv11_sunny_day_curve}
\end{figure}

Από τα επτά ΦΒ του δικτύου που είναι σηματοδοτημένα με κόκκινο, τα πέντε εμφανίζουν συντελεστή ισχύος χαμηλότερο της μονάδας κατά την πιο ηλιόλουστη ημέρα. Επιπλέον, όλα τα ΦΒ εμφανίζουν μέγιστη ημερήσια παραγωγή την ίδια ημερομηνία (2024-05-06).

\begin{figure}[H]
\centering

\begin{subfigure}[t]{0.49\textwidth}
\centering
\includegraphics[width=\linewidth]{115455-40.png}
\caption{ΦΒ:115455-\textbf{40}}
\end{subfigure}
\hfill
\begin{subfigure}[t]{0.49\textwidth}
\centering
\includegraphics[width=\linewidth]{115456-43.png}
\caption{ΦΒ:115456-\textbf{43}}
\end{subfigure}

\vspace{0.4cm}

\begin{subfigure}[t]{0.49\textwidth}
\centering
\includegraphics[width=\linewidth]{115457-46.png}
\caption{ΦΒ:115457-\textbf{46}}
\end{subfigure}
\hfill
\begin{subfigure}[t]{0.49\textwidth}
\centering
\includegraphics[width=\linewidth]{114875-19_ΠΙΟ ΗΛΙΟΛΟΥΣΤΗ ΗΜΕΡΑ ΠΑΡΑΓΩΓΗ.png}
\caption{ΦΒ:114875-\textbf{19}}
\end{subfigure}

\caption{Καμπύλες ενεργού ισχύος κατά την πιο ηλιόλουστη ημέρα για τα υπόλοιπα ΦΒ που ακολουθούν σιγμοειδή κατανομή του συντελεστή ισχύος.}
\end{figure}

 Τα παραπάνω διαγράμματα επιβεβαιώνουν ότι η συγκεκριμένη ημέρα χαρακτηρίζεται από καθαρό ουρανό, χωρίς νεφοκάλυψη, με αποτέλεσμα τα ΦΒ να λειτουργούν κοντά στη μέγιστη απόδοσή τους. Επιπλέον, ο χρόνος εμφάνισης της μέγιστης ισχύος παρουσιάζει μικρή διασπορά μεταξύ των μονάδων και εντοπίζεται στο διάστημα 12:45--13:15. Τέλος, η μορφή των καμπυλών είναι συμβατή με την αναμενόμενη μονοκορυφική, σχεδόν κανονική ημερήσια κατανομή παραγωγής.


\subsection{Διερεύνηση Γραμμικής συσχέτισης}

Στις στατιστικές, η εξάρτηση είναι οποιαδήποτε στατιστική σχέση μεταξύ δύο τυχαίων μεταβλητών ή δύο σύνολα δεδομένων\\

Για να διερεύνησουμε τη  συσχέτιση $P-\cos\phi$  θα χρησιμοποιησουμε  διάφορους συντελεστές συσχέτισης, οπού συχνά συμβολίζονται με ρ ή \textlatin{r}, μετρώντας το βαθμό συσχέτισης.
\begin{itemize}
\item Οι πιο κοινοί από αυτούς είναι ο συντελεστής συσχέτισης \textbf{\textlatin{Pearson}}, ο οποίος είναι ευαίσθητος μόνο σε μια \textbf{γραμμική} σχέση μεταξύ των δύο μεταβλητών. Ενώ άλλοι συντελεστές συσχέτισης έχουν αναπτυχθεί για να είναι πιο γεροί από το συντελεστή συσχέτισης του \textlatin{Pearson} -- που είναι πιο ευαίσθητος σε μη γραμμικές σχέσεις.

Ο συντελεστής συσχέτισης \textlatin{Pearson} ορίζεται ως:
\begin{equation}
\rho_{X,Y}=\mathrm{corr}(X,Y)=\frac{\mathrm{cov}(X,Y)}{\sigma_X\sigma_Y}
=\frac{E\!\left[(X-\mu_X)(Y-\mu_Y)\right]}{\sigma_X\sigma_Y},
\end{equation}
Οπού $\sigma_X$ και $\mathbb{E}[X]$ η τυπική απόκλιση και η μέση τιμή του $X$, αντίστοιχα.\par Η συσχέτιση \textlatin{Pearson} είναι +1 σε περίπτωση μίας τέλειας (αύξουσας) γραμμικής σχέσης (συσχέτιση), -1 σε περίπτωση μίας τέλειας φθίνουσας (αντίστροφης) γραμμικής σχέσης (αντισυσχέτιση), και κάποια τιμή μεταξύ -1 και 1 σε όλες τις άλλες περιπτώσεις, που δείχνει το βαθμό της γραμμικής εξάρτησης μεταξύ των μεταβλητών. Καθώς πλησιάζει το μηδέν υπάρχει λιγότερη  σχέση (πιο κοντά σε ασυσχέτιστα). Όσο πιο κοντά είναι ο συντελεστής είτε στο -1 ή στο 1, τόσο ισχυρότερη είναι η συσχέτιση μεταξύ των μεταβλητών
(πηγή:~\cite{wikipedia-correlation-dependence})

\item Πέρα από  τον συντελεστή \textlatin{Pearson}, εφαρμόζεται και \textbf{γραμμική παλινδρόμηση} με δείκτη αξιολόγησης τον συντελεστή παλινδρόμησης $\mathbf{R^2}$. Η χρήση του $R^2$ επιτρέπει να εκτιμηθεί πόσο καλά εξηγεί η γραμμή παλινδρόμησης τα δεδομένα κάθε ζεύγους μεταβλητών και συνεπώς αν η γραμμική προσέγγιση είναι επαρκής.
\end{itemize}

\begin{center}
\textbf{ΝΑ ΜΠΟΥΝ ΚΑΙ ΑΛΛΑ ΘΕΩΡΗΤΙΚΑ \textlatin{p-value} καμια πολυωνιμικη αμα χρειαστει\\ ...}
\end{center}

\newpage Στο πλαίσιο αυτό, εκτελέστηκε γραμμική παλινδρόμηση για τα πέντε ΦΒ που εξετάζονται, στο ζεύγος μεταβλητών $P$--$\cos\phi$. Για κάθε μονάδα υπολογίστηκαν ο συντελεστής συσχέτισης $r$ και  η ευθεία παλινδρόμησης, ώστε να αξιολογηθεί ποσοτικά ο βαθμός γραμμικής εξάρτησης.

\begin{figure}[H]
  \centering
  \includegraphics[width=0.95\textwidth]{PV_ΓΡΑΜΙΚΗ ΣΥΣΧΕΤΙΣΗ_ΑΡΧΙΚΑ_P_cosφ.png}
  \caption{Αποτελέσματα γραμμικής παλινδρόμησης για τα 5 ΦΒ στο ζεύγος $P$--$\cos\phi$ (υπολογισμός $r$ \textlatin{Pearson}).}
  \label{fig:linear_regression_p_cosphi_5pv}
\end{figure}

Γίνεται αντιληπτό από το διάγραμμα ότι δεν προκύπτει άμεσα ισχυρή γραμμική συσχέτιση, καθώς οι τιμές του $r$ είναι χαμηλές.Η χαμηλή αυτή τιμή δεν υποδηλώνει απουσία συσχέτισης, αλλά αδυναμία του γραμμικού
μοντέλου να αποτυπώσει τη σχέση στο σύνολο των δεδομένων. Ωστόσο, για μεγαλύτερες τιμές ενεργού ισχύος διαφαίνεται μια ήπια αρνητική τάση, η οποία επιβεβαιώνεται από το αρνητικό πρόσημο του συντελεστή συσχέτισης $r$.
Για να βελτιωθεί η ανάλυση, η διερεύνηση θα γίνει ξεχωριστά για κάθε ΦΒ και στη συνέχεια θα παρουσιαστούν οι αντίστοιχες παρατηρήσεις.
 
 \subsection{Προσθήκη Κατωφλιου Ισχύος για μελέτη γραμμικής συμπεριφοράς στις μεγαλύτερες τιμές}
Όπως φαίνεται στο Σχήμα~\ref{fig:linear_regression_p_cosphi_5pv}, για χαμηλές τιμές ενεργού ισχύος ο συντελεστής ισχύος εμφανίζει ασταθή συμπεριφορά, η οποία υποβαθμίζει σημαντικά τη γραμμικότητα και δυσχεραίνει την εξαγωγή αξιόπιστης συσχέτισης. Για τον λόγο αυτό, εφαρμόζεται κατώφλι ισχύος ίσο με περίπου $10\%$ της ονομαστικής ισχύος κάθε ΦΒ συστήματος.

Ως κατώφλι φιλτραρίσματος
ορίστηκαν τα $100\,\textlatin{kW}$ για τα συστήματα ΦΒ~19, ΦΒ~46 και ΦΒ~11, και τα
$50\,\textlatin{kW}$ για τα ΦΒ~40 και ΦΒ~43, λαμβάνοντας υπόψη τη μικρότερη ονομαστική
ισχύ των τελευταίων. Οι μετρήσεις κάτω του κατωφλίου αποκλείστηκαν από την ανάλυση
παλινδρόμησης λόγω της αυξημένης διακύμανσης του $\cos\varphi$ που παρατηρείται
σε συνθήκες χαμηλής παραγωγής, κατά τις οποίες ο μετατροπέας \textlatin{(inverter)} λειτουργεί εγχέοντας μεγάλα ποσοστά άεργου ισχύος ως προς την ενεργη ισχύ,αλλά χωρίς να έχει σημαντική επίδραση στις απώλειες και στη διακύμανση της τάσης του δικτύου. 

\begin{figure}
\centering
\begin{subfigure}[t]{0.49\textwidth}
\centering
\includegraphics[width=\linewidth]{linear_ΦΒ_11.png}
\caption{ΦΒ 116310-11}
\end{subfigure}
\hfill
\begin{subfigure}[t]{0.49\textwidth}
\centering
\includegraphics[width=\linewidth]{linear_ΦΒ_19.png}
\caption{ΦΒ 114875-19}
\end{subfigure}

\vspace{0.4cm}

\begin{subfigure}[t]{0.49\textwidth}
\centering
\includegraphics[width=\linewidth]{linear_ΦΒ_40.png}
\caption{ΦΒ 115455-40}
\end{subfigure}
\hfill
\begin{subfigure}[t]{0.49\textwidth}
\centering
\includegraphics[width=\linewidth]{linear_ΦΒ_43.png}
\caption{ΦΒ 115456-43}
\end{subfigure}

\vspace{0.4cm}

\begin{subfigure}[t]{0.60\textwidth}
\centering
\includegraphics[width=\linewidth]{linear_ΦΒ_46.png}
\caption{ΦΒ 115457-46}
\end{subfigure}

\caption{Κυμματομορφές μετά την εφαρμογή κατωφλίου ($\approx 10\%$ της ονομαστικής ισχύος) για τα 5 ΦΒ.}
\label{fig:linear_threshold_10pct_5pv}
\end{figure}

\newpage

Παρατηρείται ότι, μετά την εφαρμογή του κατωφλίου, ο συντελεστής συσχέτισης παρουσιάζει αισθητή βελτίωση (ενδεικτικά, από τιμές της τάξης του $r\approx 0.1$ σε τιμές περίπου $r\approx 0.5$), γεγονός που υποδηλώνει καθαρότερη γραμμική τάση στο εξεταζόμενο εύρος λειτουργίας. Ωστόσο, η γραμμική συσχέτιση παραμένει μέτρια και δεν προσεγγίζει την επιθυμητή ισχυρή συσχέτιση (τιμές κοντά στη μονάδα). Συνεπώς, απαιτείται περαιτέρω διερεύνηση εναλλακτικών προσεγγίσεων συσχέτισης. Σε πρώτο στάδιο, η ανάλυση συνεχίζεται εντός του πλαισίου της γραμμικής συσχέτισης, με την ακόλουθη παρατήρηση να αναλύεται στο επόμενο σημείο.
 
      %ΕΔΩ ΒΑΛΕ ΥΠΟΕΝΟΤΗΤΑ ΓΙΑ ΦΑΣΕΙΣ%
\subsection{Διαχωρισμός του συντελεστή ισχύος σε φάση ανόδου και καθόδου της ενεργού ισχύος}      

Παρατηρείται ότι, για ενδιάμεσες τιμές της ενεργού ισχύος $P$, εμφανίζονται σημαντικές διαφοροποιήσεις στον συντελεστή ισχύος, ακόμη και για παρόμοια επίπεδα ισχύος. Ειδικότερα, διακρίνεται εναλλασσόμενο μοτίβο τιμών (υψηλή--χαμηλή) του συντελεστή ισχύος στο ενδιάμεσο εύρος λειτουργίας.\par

%ισως αυτη η κυμματομορφη με της ωρες πρεπει να μπει πιο κατω%
\begin{figure}[!b]
  \centering
  \includegraphics[width=0.78\textwidth]{PICS/ΦΒ_11_zoom_5x5.png}
  \caption{Ενδεικτική μεγέθυνση (\textlatin{5x5}) για ΦΒ~11: στο ενδιάμεσο εύρος τιμών της ενεργού ισχύος, τα σημεία της ανερχόμενης (πρωινής) και κατερχόμενης (μεσημβρινής/απογευματινής) φάσης εμφανίζονται σε διαφορετικές στάθμες $\cos\varphi$, επιβεβαιώνοντας τη χρονική αλληλουχία που περιγράφεται στο κείμενο.}
  \label{fig:pv11_zoom_5x5_sequence}
\end{figure}

 Με βάση την ανωτέρω παρατήρηση, οι μετρήσεις κάθε ΦΒ
διαχωρίστηκαν σε δύο φάσεις: την \textit{ανερχόμενη φάση παραγωγής}, η οποία
αντιστοιχεί στο χρονικό διάστημα από την έναρξη της παραγωγής έως την επίτευξη
της μέγιστης ισχύος, και την \textit{κατερχόμενη φάση}, που καλύπτει το
διάστημα από το σημείο μέγιστης παραγωγής έως τη δύση.

%%%%%%



Παρακάτω παρατείθονται τα διαγραμματα με τη γραμμικη παλινδρόμηση χωρίζοντας τις 2 περιοχές σε ανερχόμενη και κατερχόμενη:
\begin{figure}[H]
\centering

\begin{subfigure}[t]{0.49\textwidth}
\centering
\includegraphics[width=\linewidth]{PICS/linear_phase_ΦΒ_11.png}
\caption{ΦΒ 116310-11}
\end{subfigure}
\hfill
\begin{subfigure}[t]{0.49\textwidth}
\centering
\includegraphics[width=\linewidth]{linear_phase_ΦΒ_19.png}
\caption{ΦΒ 114875-19}
\end{subfigure}

\vspace{0.4cm}

\begin{subfigure}[t]{0.49\textwidth}
\centering
\includegraphics[width=\linewidth]{linear_phase_ΦΒ_40.png}
\caption{ΦΒ 115455-40}
\end{subfigure}
\hfill
\begin{subfigure}[t]{0.49\textwidth}
\centering
\includegraphics[width=\linewidth]{linear_phase_ΦΒ_43.png}
\caption{ΦΒ 115456-43}
\end{subfigure}

\vspace{0.4cm}

\begin{subfigure}[t]{0.60\textwidth}
\centering
\includegraphics[width=\linewidth]{linear_phase_ΦΒ_46.png}
\caption{ΦΒ 115457-46}
\end{subfigure}

\caption{Κυμματομορφές διαχωριζόμενες σε πρωινές και απογευματινές ώρες για τα 5 ΦΒ.}
\label{fig:proi_apogeuma_grammiki_palidromisi}
\end{figure}


\par % ============================================================
% Ενότητα: Ανάλυση Γραμμικής Παλινδρόμησης P – cos φ ανά Φάση
% ============================================================

Τα αποτελέσματα της γραμμικής παλινδρόμησης ανά φάση παρουσιάζονται στον
Πίνακα~\ref{tab:pearson_phase}. Κατά την ανερχόμενη φάση, όλα τα συστήματα
εμφανίζουν ισχυρή αρνητική γραμμική συσχέτιση, με τιμές $r$ από $-0{,}78$
έως $-0{,}88$. Αξιοσημείωτη είναι η περίπτωση του ΦΒ~11, όπου ο συντελεστής
Pearson φτάνει $r = -0{,}88$ κατά την ανερχόμενη φάση, ενώ κατά την κατερχόμενη
φτάνει $r = -0{,}95$, αποτελώντας την ισχυρότερη γραμμική συσχέτιση που
παρατηρήθηκε στο σύνολο των εξεταζόμενων συστημάτων. Αντίθετα, τα συστήματα
ΦΒ~40 και ΦΒ~43 εμφανίζουν ισχυρή συσχέτιση μόνο κατά την ανερχόμενη φάση
($r = -0{,}78$ και $r = -0{,}85$ αντίστοιχα), ενώ κατά την κατερχόμενη φάση η
συσχέτιση είναι στατιστικά μη σημαντική ($p > 0{,}05$), υποδηλώνοντας πιθανή
λειτουργία σταθερού συντελεστή ισχύος από τον μετατροπέα κατά τις απογευματινές
ώρες.

\begin{table}[htbp]
  \centering
  \caption{Αποτελέσματα γραμμικής παλινδρόμησης $P$ -- $\cos\varphi$ ανά φάση
           παραγωγής και ΦΒ σύστημα.}
  \label{tab:pearson_phase}
  \renewcommand{\arraystretch}{1.3}
  \begin{tabular}{llcccl}
    \hline
    \textbf{ΦΒ Σύστημα} & \textbf{Φάση} & \textbf{$n$} &
    \textbf{\textlatin{Pearson} $r$} & \textbf{$p$-\textlatin{value}} & \textbf{Αξιολόγηση} \\
    \hline
    \multirow{2}{*}{ΦΒ 19}
      & Ανερχόμενη $\uparrow$ & 22 & $-0{,}870$ & $1{,}46\times10^{-7}$ & Ισχυρή \\
      & Κατερχόμενη $\downarrow$ & 23 & $-0{,}847$ & $3{,}37\times10^{-7}$ & Ισχυρή \\
    \hline
    \multirow{2}{*}{ΦΒ 40}
      & Ανερχόμενη $\uparrow$ & 21 & $-0{,}783$ & $2{,}75\times10^{-5}$ & Ισχυρή \\
      & Κατερχόμενη $\downarrow$ & 22 & $+0{,}078$ & $7{,}31\times10^{-1}$ & Μη σημαντική \\
    \hline
    \multirow{2}{*}{ΦΒ 43}
      & Ανερχόμενη $\uparrow$ & 22 & $-0{,}851$ & $5{,}30\times10^{-7}$ & Ισχυρή \\
      & Κατερχόμενη $\downarrow$ & 23 & $-0{,}186$ & $3{,}95\times10^{-1}$ & Μη σημαντική \\
    \hline
    \multirow{2}{*}{ΦΒ 46}
      & Ανερχόμενη $\uparrow$ & 21 & $-0{,}816$ & $6{,}51\times10^{-6}$ & Ισχυρή \\
      & Κατερχόμενη $\downarrow$ & 23 & $-0{,}369$ & $8{,}29\times10^{-2}$ & Οριακή \\
    \hline
    \multirow{2}{*}{ΦΒ 11}
      & Ανερχόμενη $\uparrow$ & 22 & $-0{,}885$ & $4{,}69\times10^{-8}$ & Ισχυρή \\
      & Κατερχόμενη $\downarrow$ & 24 & $-0{,}949$ & $1{,}58\times10^{-12}$ & Πολύ ισχυρή \\
    \hline
  \end{tabular}
\end{table}

Τα ευρήματα αυτά επιβεβαιώνουν ότι η χαμηλή τιμή του συντελεστή \textlatin{Pearson} επί του
συνολικού δείγματος δεν υποδηλώνει απουσία γραμμικής σχέσης, αλλά οφείλεται στην
υπέρθεση δύο διακριτών γραμμικών τάσεων με διαφορετικά χαρακτηριστικά. Κατά την
ανερχόμενη φάση, ο συντελεστής $\cos\varphi$ κινείται κοντά στη μονάδα και μειώνεται
σταδιακά καθώς αυξάνεται η παραγωγή, αντανακλώντας την προοδευτική ενεργοποίηση
του ελέγχου αέργου ισχύος από τον μετατροπέα. Κατά την κατερχόμενη φάση, για ίδιο
επίπεδο ισχύος $P$, ο $\cos\varphi$ παρουσιάζει χαμηλότερες τιμές σε σχέση με
την ανερχόμενη φάση, γεγονός που αποδίδεται στη συνεχή παροχή αέργου ισχύος στο
δίκτυο κατά τις μεσημβρινές ώρες 
από τον μετατροπέα και κατά την πτωτική πορεία της παραγωγής. Το φαινόμενο αυτό εξηγεί γιατί η
χρήση μόνο της ενεργού ισχύος ως ανεξάρτητης μεταβλητής δεν επαρκεί για την πλήρη
περιγραφή της συμπεριφοράς του συντελεστή ισχύος σε φωτοβολταϊκά συστήματα
συνδεδεμένα στο δίκτυο.\\
\textbf{Ετσι σε επόμενο επίπεδο θα μελετηθεί  η επίδραση της αεργου ισχύος και της τάσης \textlatin{V-cos}φ ,\textlatin{Q-cos}φ}.

\begin{figure}[H]
  \centering
  \includegraphics[width=1\textwidth]{PICS/PV_linear_vs_poly.png}
  \caption{Σύγκριση γραμμικής και πολυωνυμικής προσαρμογής στο ζεύγος $P$--$\cos\varphi$ (\textlatin{linear vs poly}).}
  \label{fig:linear_vs_poly_p_cosphi}
\end{figure}

Η σύγκριση του Σχήματος~\ref{fig:linear_vs_poly_p_cosphi} δείχνει ότι η πολυωνυμική προσαρμογή αποδίδει καλύτερα τη σχέση μεταξύ $P$ και $\cos\varphi$ από τη γραμμική. Ο βασικός λόγος είναι ότι τα δεδομένα δεν ακολουθούν ευθεία, αλλά παρουσιάζουν καμπυλότητα και διαφορετική κλίση στα επιμέρους εύρη ισχύος. Ένα γραμμικό μοντέλο επιβάλλει σταθερή κλίση σε όλο το πεδίο και οδηγεί σε συστηματικά σφάλματα, ενώ η πολυωνυμική μορφή επιτρέπει μεταβαλλόμενη κλίση και προσαρμόζεται αποδοτικότερα στη μη γραμμική συμπεριφορά των μετρήσεων. Συνεπώς, η πολυωνυμική συσχέτιση αποτελεί καταλληλότερη περιγραφή για το συγκεκριμένο σύνολο δεδομένων και δικαιολογεί την περαιτέρω διερεύνηση μη γραμμικών μοντέλων.




\section{συγκριση \textlatin{P,Q,V-COS}φ για πιο απομακρυσμένο κόμβο στις μετρήσεις του δικτύου}

Για τον κόμβο 48 (ή τον κόμβο 117168), που αποτελεί τον πιο απομακρυσμένο κόμβο του δικτύου, λαμβάνονται μετρήσεις των μεγεθών P, Q και V κατά τη διάρκεια μιας ηλιόλουστης ημέρας. Στη συνέχεια, τα δεδομένα απεικονίζονται στο ίδιο διάγραμμα, με στόχο τη σύγκριση και την ανάλυση της επίδρασης στις μεταβολές του συντελεστή ισχύος 
\begin{figure}[htbp]
  \centering
  \begin{tikzpicture}
    \begin{axis}[
      width=0.88\textwidth,
      height=7cm,
      axis y line*=left,
      axis x line*=bottom,
      xlabel={},
      ylabel={$\cos\varphi$},
      ylabel style={color=blue!60!black},
      ymin=0.989783, ymax=1.000929,
      scaled y ticks=false,
      yticklabel style={{color={cf_color}, font=\small,
        /pgf/number format/fixed,
        /pgf/number format/precision=4}},
      xmin=-0.5, xmax=54.5,
      xtick={2,6,10,14,18,22,26,30,34,38,42,46,50,54},
      xticklabels={07:00,08:00,09:00,10:00,11:00,12:00,13:00,14:00,15:00,16:00,17:00,18:00,19:00,20:00},
      x tick label style={rotate=45, anchor=east, font=\small},
      y tick label style={color=blue!60!black, font=\small},
      tick align=outside,
      ymajorgrids=true,
      grid style={dotted, gray!40},
      legend style={at={(0.02,0.98)}, anchor=north west, font=\small,
                    draw=gray!50, fill=white, fill opacity=0.85},
    ]
      \addplot[color=blue!60!black, thick, mark=*, mark size=1.2pt]
        coordinates { (0,0.990712) (1,0.998488) (2,0.999392) (3,0.999945) (4,0.999981) (5,0.999997) (6,0.999973) (7,0.999931) (8,0.999890) (9,0.999849) (10,0.999797) (11,0.999673) (12,0.999587) (13,0.999516) (14,0.999464) (15,0.999340) (16,0.999211) (17,0.999154) (18,0.999130) (19,0.999046) (20,0.999001) (21,0.998935) (22,0.999024) (23,0.998919) (24,0.998875) (25,0.998883) (26,0.998923) (27,0.998898) (28,0.998977) (29,0.998945) (30,0.999091) (31,0.999000) (32,0.998990) (33,0.999024) (34,0.999085) (35,0.999127) (36,0.999226) (37,0.999270) (38,0.999318) (39,0.999395) (40,0.999483) (41,0.999543) (42,0.999596) (43,0.999697) (44,0.999796) (45,0.999858) (46,0.999846) (47,0.999950) (48,0.999986) (49,1.000000) (50,1.000000) (51,0.999904) (52,0.999973) (53,0.999929) (54,0.997785) };
      \addlegendentry{$\cos\varphi$}
    \end{axis}

    \begin{axis}[
      width=0.92\textwidth,
      height=7cm,
      axis y line*=right,
      axis x line=none,
      ylabel={$P$ \textlatin{(MW)}},
      ylabel style={color=teal!70!black},
      ymin=-0.073498, ymax=0.723178,
      scaled y ticks=false,
      xmin=-0.5, xmax=54.5,
      y tick label style={color=teal!70!black, font=\small,
        /pgf/number format/fixed,
        /pgf/number format/precision=3},
      legend style={at={(0.98,0.98)}, anchor=north east, font=\small,
                    draw=gray!50, fill=white, fill opacity=0.85},
    ]
      \addplot[color=teal!70!black, thick, dashed, mark=square*, mark size=1.2pt]
        coordinates { (0,0.006120) (1,0.013080) (2,0.020640) (3,0.034440) (4,0.058440) (5,0.091560) (6,0.131040) (7,0.173640) (8,0.218640) (9,0.262200) (10,0.303480) (11,0.342720) (12,0.379680) (13,0.416400) (14,0.450480) (15,0.485400) (16,0.516240) (17,0.542400) (18,0.566280) (19,0.584880) (20,0.600840) (21,0.613200) (22,0.624360) (23,0.631920) (24,0.639600) (25,0.644280) (26,0.645960) (27,0.646080) (28,0.644040) (29,0.636960) (30,0.621720) (31,0.630240) (32,0.618960) (33,0.607800) (34,0.594240) (35,0.576960) (36,0.554760) (37,0.530640) (38,0.503400) (39,0.472320) (40,0.440040) (41,0.404880) (42,0.367080) (43,0.326760) (44,0.284880) (45,0.241920) (46,0.198480) (47,0.155520) (48,0.113520) (49,0.075000) (50,0.045000) (51,0.026040) (52,0.016200) (53,0.010080) (54,0.003600) };
      \addlegendentry{$P$ (\textlatin{MW})}
    \end{axis}
  \end{tikzpicture}
  \caption{Συσχέτιση $\cos\varphi$ και ενεργού ισχύος $P$}
  \label{fig:p_cosphi}
\end{figure}

\begin{figure}[htbp]
  \centering
  \begin{tikzpicture}
    \begin{axis}[
      width=0.92\textwidth,
      height=7cm,
      axis y line*=left,
      axis x line*=bottom,
      xlabel={},
      ylabel={$\cos\varphi$},
      ylabel style={color=blue!60!black},
      ymin=0.989783, ymax=1.000929,
      scaled y ticks=false,
      yticklabel style={{color={cf_color}, font=\small,
       /pgf/number format/fixed,
       /pgf/number format/precision=4}},
      xmin=-0.5, xmax=54.5,
      xtick={2,6,10,14,18,22,26,30,34,38,42,46,50,54},
      xticklabels={07:00,08:00,09:00,10:00,11:00,12:00,13:00,14:00,15:00,16:00,17:00,18:00,19:00,20:00},
      x tick label style={rotate=45, anchor=east, font=\small},
      y tick label style={color=blue!60!black, font=\small},
      tick align=outside,
      ymajorgrids=true,
      grid style={dotted, gray!40},
      legend style={at={(0.02,0.98)}, anchor=north west, font=\small,
                    draw=gray!50, fill=white, fill opacity=0.85},
    ]
      \addplot[color=blue!60!black, thick, mark=*, mark size=1.2pt]
        coordinates { (0,0.990712) (1,0.998488) (2,0.999392) (3,0.999945) (4,0.999981) (5,0.999997) (6,0.999973) (7,0.999931) (8,0.999890) (9,0.999849) (10,0.999797) (11,0.999673) (12,0.999587) (13,0.999516) (14,0.999464) (15,0.999340) (16,0.999211) (17,0.999154) (18,0.999130) (19,0.999046) (20,0.999001) (21,0.998935) (22,0.999024) (23,0.998919) (24,0.998875) (25,0.998883) (26,0.998923) (27,0.998898) (28,0.998977) (29,0.998945) (30,0.999091) (31,0.999000) (32,0.998990) (33,0.999024) (34,0.999085) (35,0.999127) (36,0.999226) (37,0.999270) (38,0.999318) (39,0.999395) (40,0.999483) (41,0.999543) (42,0.999596) (43,0.999697) (44,0.999796) (45,0.999858) (46,0.999846) (47,0.999950) (48,0.999986) (49,1.000000) (50,1.000000) (51,0.999904) (52,0.999973) (53,0.999929) (54,0.997785) };
      \addlegendentry{$\cos\varphi$}
    \end{axis}

    \begin{axis}[
      width=0.92\textwidth,
      height=7cm,
      axis y line*=right,
      axis x line=none,
      ylabel={$Q$ \textlatin{(MVAr)}},
      ylabel style={color=orange!80!black},
      ymin=-0.004598, ymax=0.034238,
      scaled y ticks=false,
      xmin=-0.5, xmax=54.5,
      y tick label style={color=orange!80!black, font=\small,
        /pgf/number format/fixed,
        /pgf/number format/precision=4},
      legend style={at={(0.98,0.98)}, anchor=north east, font=\small,
                    draw=gray!50, fill=white, fill opacity=0.85},
    ]
      \addplot[color=orange!80!black, thick, dashed, mark=square*, mark size=1.2pt]
        coordinates { (0,-0.000840) (1,-0.000720) (2,-0.000720) (3,-0.000360) (4,-0.000360) (5,0.000240) (6,0.000960) (7,0.002040) (8,0.003240) (9,0.004560) (10,0.006120) (11,0.008760) (12,0.010920) (13,0.012960) (14,0.014760) (15,0.017640) (16,0.020520) (17,0.022320) (18,0.023640) (19,0.025560) (20,0.026880) (21,0.028320) (22,0.027600) (23,0.029400) (24,0.030360) (25,0.030480) (26,0.030000) (27,0.030360) (28,0.029160) (29,0.029280) (30,0.026520) (31,0.028200) (32,0.027840) (33,0.026880) (34,0.025440) (35,0.024120) (36,0.021840) (37,0.020280) (38,0.018600) (39,0.016440) (40,0.014160) (41,0.012240) (42,0.010440) (43,0.008040) (44,0.005760) (45,0.004080) (46,0.003480) (47,0.001560) (48,0.000600) (49,0.000000) (50,0.000000) (51,-0.000360) (52,-0.000120) (53,-0.000120) (54,-0.000240) };
      \addlegendentry{$Q$ \textlatin{(MVAr)}}
    \end{axis}
  \end{tikzpicture}
  \caption{Συσχέτιση $\cos\varphi$ και αέργου ισχύος $Q$}
  \label{fig:q_cosphi}
\end{figure}

\begin{figure}[htbp]
  \centering
  \begin{tikzpicture}
    \begin{axis}[
      width=0.92\textwidth,
      height=7cm,
      axis y line*=left,
      axis x line*=bottom,
      xlabel={},
      ylabel={$\cos\varphi$},
      ylabel style={color=blue!60!black},
      ymin=0.989783, ymax=1.000929,
      scaled y ticks=false,
      yticklabel style={{color={cf_color}, font=\small,
       /pgf/number format/fixed,
       /pgf/number format/precision=4}},
      xmin=-0.5, xmax=54.5,
      xtick={2,6,10,14,18,22,26,30,34,38,42,46,50,54},
      xticklabels={07:00,08:00,09:00,10:00,11:00,12:00,13:00,14:00,15:00,16:00,17:00,18:00,19:00,20:00},
      x tick label style={rotate=45, anchor=east, font=\small},
      y tick label style={color=blue!60!black, font=\small},
      tick align=outside,
      ymajorgrids=true,
      grid style={dotted, gray!40},
      legend style={at={(0.02,0.98)}, anchor=north west, font=\small,
                    draw=gray!50, fill=white, fill opacity=0.85},
    ]
      \addplot[color=blue!60!black, thick, mark=*, mark size=1.2pt]
        coordinates { (0,0.990712) (1,0.998488) (2,0.999392) (3,0.999945) (4,0.999981) (5,0.999997) (6,0.999973) (7,0.999931) (8,0.999890) (9,0.999849) (10,0.999797) (11,0.999673) (12,0.999587) (13,0.999516) (14,0.999464) (15,0.999340) (16,0.999211) (17,0.999154) (18,0.999130) (19,0.999046) (20,0.999001) (21,0.998935) (22,0.999024) (23,0.998919) (24,0.998875) (25,0.998883) (26,0.998923) (27,0.998898) (28,0.998977) (29,0.998945) (30,0.999091) (31,0.999000) (32,0.998990) (33,0.999024) (34,0.999085) (35,0.999127) (36,0.999226) (37,0.999270) (38,0.999318) (39,0.999395) (40,0.999483) (41,0.999543) (42,0.999596) (43,0.999697) (44,0.999796) (45,0.999858) (46,0.999846) (47,0.999950) (48,0.999986) (49,1.000000) (50,1.000000) (51,0.999904) (52,0.999973) (53,0.999929) (54,0.997785) };
      \addlegendentry{$\cos\varphi$}
    \end{axis}

    \begin{axis}[
      width=0.92\textwidth,
      height=7cm,
      axis y line*=right,
      axis x line=none,
      ylabel={$V$ (\textlatin{pu})},
      ylabel style={color=violet!70!black},
      ymin=0.987648, ymax=1.119412,
      scaled y ticks=false,
      xmin=-0.5, xmax=54.5,
      y tick label style={color=violet!70!black, font=\small,
        /pgf/number format/fixed,
        /pgf/number format/precision=3},
      legend style={at={(0.98,0.98)}, anchor=north east, font=\small,
                    draw=gray!50, fill=white, fill opacity=0.85},
    ]
      \addplot[color=violet!70!black, thick, dashed, mark=square*, mark size=1.2pt]
        coordinates { (0,1.000920) (1,1.002667) (2,1.004759) (3,1.007732) (4,1.012159) (5,1.017869) (6,1.024067) (7,1.030416) (8,1.036713) (9,1.042759) (10,1.048434) (11,1.053851) (12,1.058799) (13,1.066440) (14,1.079746) (15,1.084829) (16,1.089374) (17,1.093184) (18,1.096348) (19,1.099025) (20,1.101340) (21,1.103247) (22,1.104254) (23,1.105489) (24,1.106222) (25,1.106461) (26,1.106586) (27,1.106661) (28,1.106361) (29,1.106110) (30,1.105255) (31,1.104774) (32,1.103804) (33,1.102369) (34,1.100253) (35,1.098091) (36,1.095173) (37,1.091368) (38,1.087686) (39,1.082728) (40,1.077242) (41,1.071714) (42,1.066089) (43,1.059091) (44,1.052079) (45,1.044881) (46,1.037547) (47,1.029807) (48,1.022214) (49,1.015192) (50,1.009565) (51,1.005804) (52,1.003325) (53,1.001806) (54,1.000399) };
      \addlegendentry{$V$ (\textlatin{pu})}
    \end{axis}
  \end{tikzpicture}
  \caption{Συσχέτιση $\cos\varphi$ και τάσης $V$}
  \label{fig:v_cosphi}
\end{figure}

\newpage
Από την ανάλυση των διαγραμμάτων προκύπτει ότι κατά τις μεσημβρινές ώρες, η αυξημένη έγχυση ενεργού ισχύος (\textlatin{P}) από τα φωτοβολταϊκά συνοδεύεται από άνοδο της τάσης ({\textlatin{V})(πέρα των επιτρεπτών ορίων $>$ 1.1 \textlatin{p.u}) στον κόμβο. Παράλληλα, παρατηρείται αύξηση της άεργου ισχύος (\textlatin{Q}), γεγονός που οδηγεί σε μείωση του συντελεστή ισχύος. Συνεπώς, διαπιστώνεται ότι η ταυτόχρονη αύξηση των μεγεθών \textlatin{P, V} και \textlatin{Q} συσχετίζεται με υποβάθμιση του συντελεστή ισχύος 

\section{Συσχέτιση τάσεων \textlatin{Power Factory} με μετρήσεις ΔΕΔΔΗΕ}
Στο τελευταίο στάδιο της ανάλυσης εξετάζεται η αντιστοίχιση των μετρήσεων τάσης του ΔΕΔΔΗΕ με τα αποτελέσματα προσομοίωσης του \textlatin{Power Factory} (\textlatin{PF}) για τον ζυγό \textlatin{SB2} και τον πιο απομακρυσμένο κόμβο του δικτύου. Επειδή δεν υπάρχουν διαθέσιμες μετρήσεις τάσης ΔΕΔΔΗΕ για την 2024-05-06, χρησιμοποιείται η 2024-05-22, η οποία αντιστοιχεί επίσης σε ηλιόλουστη ημέρα με συγκρίσιμο προφίλ λειτουργίας.

\begin{figure}[H]
  \centering
  \includegraphics[width=0.92\textwidth]{PICS/deddie_pf_voltage_comparison.png}
  \caption{Σύγκριση τάσεων ΔΕΔΔΗΕ και \textlatin{Power Factory} (ΠΡΙΝ ΚΑΙ ΜΕΤΑ ΤΟΝ ΜΣ)  στον απομακρυσμένο κόμβο.}
  \label{fig:deddie_pf_voltage_comparison}
\end{figure}

Όπως φαίνεται στο Σχήμα~\ref{fig:deddie_pf_voltage_comparison}, οι δύο χρονοσειρές εμφανίζουν σαφή συσχέτιση ως προς τη χρονική μεταβολή (παρόμοια ημερήσια τάση), γεγονός που υποδεικνύει ότι το μοντέλο \textlatin{PF} αναπαράγει ικανοποιητικά τη δυναμική του συστήματος. Παράλληλα, διακρίνεται μόνιμη απόκλιση μεταξύ των καμπυλών (συστηματικό \textlatin{offset}), καθώς οι υπολογισμένες τιμές από το \textlatin{Power Factory} παραμένουν σε διαφορετικό επίπεδο από τις μετρήσεις του ΔΕΔΔΗΕ σε όλη τη διάρκεια της ημέρας. Επιπλέον, το εύρος διακύμανσης των μετρήσεων ΔΕΔΔΗΕ είναι μικρότερο από το αντίστοιχο εύρος των τιμών που προκύπτουν από το \textlatin{PF}, ένδειξη ότι το μοντέλο δίνει εντονότερη μεταβολή τάσης από αυτήν που καταγράφεται στο πεδίο.
\textbf{Σημείωση:} Οι μετρήσεις τάσης του ΔΕΔΔΗΕ αφορούν μέση τιμή (\textlatin{mean}) και όχι επιμέρους φασικές τάσεις, ενώ εκφράζονται σε \textlatin{p.u}.

\subsection{Τάση υποσταθμού ΔΕΔΔΗΕ έναντι \textlatin{Power Factory}}
\begin{figure}[H]
  \centering
  \includegraphics[width=0.92\textwidth]{PICS/ΤΑΣΗ ΥΣ ΔΕΔΔΗΕ VS PF.png}
  \caption{Σύγκριση τάσης υποσταθμού: ΔΕΔΔΗΕ \textlatin{vs} \textlatin{Power Factory}.}
  \label{fig:substation_voltage_deddie_vs_pf}
\end{figure}

Στο Σχήμα~\ref{fig:substation_voltage_deddie_vs_pf} παρατηρείται ότι η τάση του \textlatin{Power Factory} εμφανίζεται σχεδόν αμετάβλητη (\textit{άκλοντη}), επειδή στο μοντέλο ο υποσταθμός έχει θεωρηθεί ως ισχυρός κόμβος αναφοράς (\textlatin{busbar}), με υψηλή ευρωστία απέναντι σε μεταβολές τάσης. Αντίθετα, οι πραγματικές μετρήσεις του ΔΕΔΔΗΕ είναι πιο ευάλωτες στις συνθήκες λειτουργίας του πεδίου (μεταβολές φορτίου, τοπικές διαταραχές και λειτουργικά φαινόμενα του δικτύου), με αποτέλεσμα να παρουσιάζουν μεγαλύτερη διακύμανση.
Τώρα απαιτείται επόμενο βήμα βαθμονόμησης ώστε οι δύο καμπύλες να συγκλίνουν.Δηλαδή Θελω να αλλαζω το \textlatin{voltage setpoint} του \textlatin{external grid} δυναμικά για αυτο θα χρησιμοποιησω η 
\textlatin{DPL or python} με βασει της μετρησεις θελω να το κανω test.
Αυτο το θελουμε για να μην εχουμε αποκλιση στους αλλους ζυγους του δικτυου που οφειλεται στη θεωρηση σταθερης τασης στο \textlatin{slack bus}.Απο το \textlatin{shift} τασης μεταξυ μετρησεων και \textlatin{load flow} που θα προκυψει από τον απομακρυσμένο ζυγό θα ξέρουμε ενδεικτικά πως θα αλλάξουμε την τάση στο \textlatin{slack bus}.\textcolor{blue}{ΜΗΤΣΟ ΒΑΛΕ ΤΗΝ ΙΔΙΑ ΗΜΕΡΟΜΗΝΙΑ ΜΕ ΤΟΝ ΑΠΟΜΑΚΡΥΣΜΕΝΟ 22/5/2024 ΟΧΙ ΑΛΛΗ ΓΙΑ ΝΑ ΜΠΟΡΕΙΣ ΝΑ ΔΙΚΑΙΟΛΟΓΗΣΕΙΣ ΟΡΘΑ ΤΙΣ ΑΠΟΚΛΙΣΕΙΣ,ΠΑΙΞΕ ΔΗΛΑΔΗ ΜΕ ΑΥΤΗ ΤΗΝ ΗΜΕΡΟΜΝΙΑ ΜΟΝΟ}.


\subsection{ΕΛΕΓΧΟΣ ΤΑΣΗΣ ΜΕ \textlatin{VOLTAGE SETPOINT}}
Για τη μείωση της απόκλισης μεταξύ προσομοίωσης και μετρήσεων, απαιτείται  στον ζυγό 0, δηλαδή το \textlatin{voltage setpoint} στη χαμηλη πλευρα του ΜΣ που ενωνεται με το δίκτυο (\textlatin{PCC)}),να εχει δυναμική μεταβολή της τάσης ανά 15 λεπτά, σύμφωνα με το χρονικό προφίλ λειτουργίας. Οι αποκλίσεις που παρατηρούνται οφείλονται σε σημαντικό βαθμό στο ότι η τάση στον κεντρικό ζυγό έχει θεωρηθεί σταθερή, γεγονός που επηρεάζει και τους υπόλοιπους ζυγούς του μοντέλου. Σημειώνεται ότι η επιθυμητή προσέγγιση δεν είναι \textlatin{voltage control} προς μία σταθερή τιμή-στόχο, αλλά χρονομεταβαλλόμενο \textlatin{voltage setpoint}, ώστε η είσοδος της προσομοίωσης να ακολουθεί τις πραγματικές συνθήκες του δικτύου.
Για την υλοποίηση της παραπάνω στρατηγικής δημιουργήθηκε νέο \textlatin{variation} με όνομα \textlatin{ExtGrid\_LV\_Boundary}. Στο σενάριο αυτό, το \textlatin{external grid} αφαιρέθηκε από την υψηλή πλευρά και τοποθετήθηκε στη χαμηλή πλευρά του Μ/Σ, ώστε το \textlatin{setpoint} τάσης να επιβάλλεται άμεσα από το \textlatin{external grid} στον ζυγό ενδιαφέροντος.

\section{ΑΠΟΚΕΝΤΡΩΜΕΝΟΙ ΕΛΕΓΧΟΙ ΔΙΚΤΥΟΥ}
Παρατηρείται στις μετρήσεις του δικτύου μεγάλη ανύψωση τάσης τόσο στον απομακρυσμένο κόμβο(47) $δV>1.1 p.u$ όσο και στον  υποσταθμό  \textlatin{PCC} $δV>1.06 p.u$
\paragraph{Επιτρεπόμενες διακυμάνσεις τάσης}
Σύμφωνα με τους Κανονισμούς Παροχής Ηλεκτρικής Ενέργειας, η τάση σε μόνιμη κατάσταση για συστήματα με επίπεδα τάσης από $1~kV$  έως $132~kV$  άρα και στα $20kV$ οφείλει, να διατηρείται εντός του εύρους $\pm 6\%$ της ονομαστικής τιμής.~\cite{masters2002voltageRise}.
Ενώ στο \textlatin{PCC} εφαρμόζεται αυστηρότερο όριο της τάξης του $\pm 3\%$ της ονομαστικής τιμής, όπως αναφέρεται και στη σχετική βιβλιογραφία για τη διασύνδεση φωτοβολταϊκών σε \textcolor{blue}{δίκτυα χαμηλής τάσης}~\cite{stetz2012lowVoltagePV}.


Συνεπώς, για τη διατήρηση της τάσης εντός των επιτρεπόμενων ορίων στο υπό εξέταση δίκτυο, κρίνεται αναγκαία η εφαρμογή ελέγχου άεργου ισχύος (\textlatin{RPC})(\textlatin{Reactive Power Control}) στις μονάδες \textlatin{DG}. Η μέθοδος αυτή αποτελεί αποτελεσματική λύση για τον περιορισμό υπερτάσεων, υπερέχοντας έναντι τεχνικών περιορισμού ενεργού ισχύος, ιδίως λόγω του χαμηλού λόγου $\mathrm{R/X}$ των γραμμών μέσης τάσης.~\cite{kryonidis2017probabilisticVoltageControl}
Η μέθοδος αυτή βασίζεται στη \textbf{ρύθμιση της έγχυσης άεργου ισχύος από κάθε αντιστροφέα με στόχο τη σταθεροποίηση της τοπικής τάσης}. Η λειτουργία της είναι αυτόνομη και αποκεντρωμένη δηλαδή χωρίς την ανάγκη επικοινωνίας μεταξύ των μονάδων καθώς δεν απαιτείται κεντρικός συντονισμός από τον διαχειριστή του δικτύου. Αντίθετα, κάθε φωτοβολταϊκή μονάδα ανταποκρίνεται \textbf{ανεξάρτητα} στις τοπικές συνθήκες λειτουργίας, αξιοποιώντας μετρήσεις όπως η ενεργός ισχύς $P$, η άεργος ισχύς $Q$, η τάση $V$ και ο συντελεστής ισχύος $\cos\varphi$. 

\textcolor{blue}{ΝΑ ΠΡΟΣΤΕΘΕΙ ΠΕΡΙΣΣΟΤΕΡΟ ΘΕΩΡΗΤΙΚΟ ΥΠΟΒΑΘΡΟ ΓΙΑ ΤΟ ΤΙ ΕΙΝΑΙ ΤΟ \textlatin{RPC}, ΠΩΣ ΛΕΙΤΟΥΡΓΕΙ ΚΑΙ ΠΟΙΑ ΕΙΝΑΙ ΤΑ ΒΑΣΙΚΑ ΠΛΕΟΝΕΚΤΗΜΑΤΑ/ΜΕΙΟΝΕΚΤΗΜΑΤΑ ΤΟΥ.}\\
Θα μελετηθει η επίδραση της μεθόδου αυτη μέσω τριων διαφορετικών ελέγχων \textbf{\textlatin{Q(V),Q(P)},$cos$φ-\textlatin{P}}.\\

\subsection{Όρια έγχυσης άεργης ισχύος αντιστροφέα}

Οι μονάδες \textlatin{DER} πρέπει να είναι ικανές τόσο να εγχέουν άεργο ισχύ (\textlatin{over-excited}) όσο και να απορροφούν άεργο ισχύ (\textlatin{under-excited}) για επίπεδα ενεργού ισχύος μεγαλύτερα ή ίσα από την ελάχιστη μόνιμη ενεργό ισχύ \textlatin{$P_{\min}$}, ή από το \textlatin{5\%} της ονομαστικής ενεργού ισχύος \textlatin{$P_{rated}$}, όποιο από τα δύο είναι μεγαλύτερο.

Για λειτουργία με ενεργό ισχύ μεγαλύτερη από \textlatin{5\%} και μικρότερη από \textlatin{20\%} της ονομαστικής ενεργού ισχύος, η δυνατότητα ανταλλαγής άεργου ισχύος κλιμακώνεται αναλογικά με την ενεργό ισχύ και προκύπτει από τις ελάχιστες τιμές του Πίνακα~\ref{tab:der-reactive-capability}.

Για λειτουργία πάνω από το \textlatin{20\%} της ονομαστικής ενεργού ισχύος, η παροχή ή απορρόφηση άεργου ισχύος δεν πρέπει να περιορίζεται, μέχρι τα όρια ικανότητας που δίνονται στον Πίνακα~\ref{tab:der-reactive-capability}, σύμφωνα με την εκάστοτε ενεργή λειτουργία ελέγχου. Επιτρέπεται επίσης περιορισμός της ενεργού ισχύος, εφόσον απαιτείται για την τήρηση των περιορισμών φαινόμενης ισχύος.

Οι απαιτήσεις αυτές συνοψίζονται στον ακόλουθο πίνακα:

\begin{table}[htbp!]
    \centering
    \caption{Ελάχιστη ικανότητα έγχυσης και απορρόφησης άεργου ισχύος των μονάδων \textlatin{DER}.}
    \label{tab:der-reactive-capability}
    \begin{tabular}{|c|c|c|}
        \hline
        \textbf{Κατηγορία} & \textbf{Ικανότητα έγχυσης (\% της ονομαστικής φαινόμενης ισχύος)} & \textbf{Ικανότητα απορρόφησης (\% της ονομαστικής φαινόμενης ισχύος)} \\ 
        \hline
        A (στην ονομαστική τάση του \textlatin{DER}) & 44 & 25 \\ 
        \hline
        B (σε όλο το εύρος \textlatin{ANSI C84.1} περιοχή A) & 44 & 44 \\ 
        \hline
    \end{tabular}
\end{table} 

\newpage
Η μονάδα \textlatin{DER} πρέπει να μπορεί να παράγει άεργο ισχύ τουλάχιστον $44 \%$ της ονομαστικής φαινόμενης ισχύος αλλά και ενεργό ισχύ μέχρι το όριο της ονομαστικής φαινόμενης ισχύος, υπό την προϋπόθεση ότι παραμένει συνεχώς ικανή να παρέχει ή να απορροφά άεργο ισχύ σε όλο το απαιτούμενο εύρος λειτουργίας, ανάλογα με τον επιλεγμένο τρόπο ελέγχου άεργου ισχύος και τις αντίστοιχες ρυθμίσεις του διαχειριστή του δικτύου, σύμφωνα με το πρότυπο \textlatin{IEEE 1547-2018}.~\cite{ieee1547_2018}.
 
Γενικά σε χαμηλή ενεργό ισχύ:

\begin{itemize}
    \item ο αντιστροφέας έχει διαθέσιμη φαινόμενη ισχύ
    \item μπορεί να παρέχει άεργο ισχύ
    \item αυτή χρησιμοποιείται για υποστήριξη τάσης
\end{itemize}

Ενώ σε υψηλή ενεργό ισχύ:

\textbf{Υπάρχει περιορισμός της ονομαστικής ικανότητας φαινόμενης ισχύος \textlatin{(S)} του αντιστροφέα }
καθώς η ενεργός και άεργος ισχύς συνδέονται με:

\begin{equation}
S^2 = P^2 + Q^2
\end{equation}

Άρα:
\begin{itemize}
    \item Αν το ΦΒ παράγει πολλή ενεργό ισχύ ($P$), υπάρχει λίγο περιθώριο για άεργο ($Q$).
    \item Αν το ΦΒ παράγει λίγη ενεργό ισχύ, υπάρχει μεγάλο περιθώριο για άεργο ισχύ.
\end{itemize}
Σύμφωνα με τη βιβλιογραφία, σε συνθήκες αυξημένης παραγωγής, η απαίτηση για έγχυση άεργου ισχύος προς ρύθμιση της τάσης ενδέχεται να περιορίζεται από το όριο φαινόμενης ισχύος του αντιστροφέα. Σε αυτή την περίπτωση, υιοθετείται προτεραιότητα άεργου ισχύος, οδηγώντας σε αντίστοιχη περικοπή της ενεργού ισχύος. 


 
\subsection{Έλεγχος \textlatin{Q-V  Droop}}
 
Ο έλεγχος άεργου ισχύος ως συνάρτηση της τάσης, γνωστός ως \textlatin{Q(V) droop control}, αποτελεί μία ευρέως χρησιμοποιούμενη τοπική στρατηγική ρύθμισης τάσης σε δίκτυα με υψηλή διείσδυση διανεμημένης παραγωγής. Η βασική αρχή λειτουργίας του συνίσταται στη δυναμική προσαρμογή της εγχεόμενης ή απορροφούμενης άεργου ισχύος από τις μονάδες παραγωγής, σε συνάρτηση με τις τοπικές αποκλίσεις της τάσης από την ονομαστική τιμή.

Συγκεκριμένα,ο στόχος της μεθόδου είναι διττός,πρώτον όταν η τάση στον κόμβο υπερβαίνει ένα προκαθορισμένο όριο, η μονάδα απορροφά άεργη ισχύ, συμβάλλοντας στη μείωση της τάσης, ενώ σε περιπτώσεις χαμηλής τάσης εγχέει άεργη ισχύ για την υποστήριξή της.Δεύτερον να διατηρείται μηδενική άεργος ισχύς για ένα σχετικά ευρύ λειτουργικό εύρος τάσης, ώστε να αποφεύγονται περιττές παρεμβάσεις άεργου ισχύος όταν η τάση παραμένει σε αποδεκτά επίπεδα. Με αυτόν τον τρόπο αποφεύγεται άσκοπη κατανάλωση/παραγωγή ενεργού ισχύος και συνεπώς μειώνονται οι απώλειες του δικτύου. Η σχέση αυτή περιγράφεται συνήθως μέσω μιας γραμμικής χαρακτηριστικής (\textlatin{droop curve}), η οποία καθορίζεται από κατάλληλες παραμέτρους κλίσης και ορίων λειτουργίας.Ενώ η απόδοση της μεθόδου αυτής βασίζεται στο χαμηλό σχετικά λόγο $\mathrm{R/X}$ των γραμμών μέσης τάσης κάτι που δεν ισχύει αν τα ΦΒ ήταν συνδεδεμένα στη χαμηλή τάση.


Μια τυπική μορφή της χαρακτηριστικής στατισμού \textlatin{$Q=f(V)$} φαίνεται στο Σχήμα~\ref{fig:qv-kryonidis}, όπου η άεργος ισχύς της μονάδας \textlatin{DG} καθορίζεται σε συνάρτηση με την τάση στο σημείο κοινής σύνδεσης (\textlatin{PCC}).

\begin{figure}[H]
    \centering
    \includegraphics[width=0.78\textwidth]{qv-kryonidis.png}
    \caption{Χαρακτηριστική \textlatin{Q(V)} σύμφωνα με τη μεθοδολογία του Κρυωνίδη.~\cite{dimoulias2024droop}}
    \label{fig:qv-kryonidis}
\end{figure}

Η συγκεκριμένη κυμματομορφή ενσωματτώνεται στον έλεγχο \textlatin{Q(V)} για τους αντιστροφεις του δικτυου με πρωταρχικό στόχο την βελτίωση της τάσης που υπερβαίνει τα επιτρεπτά όρια με τη παρούσα λειτουργία, παρακολουθώντας παράλληλα της απώλειες αλλά και τα ποσά έγχυσης άεργου ισχύος.
Η συγκεκριμένη χαρακτηριστική καμπύλη ενσωματώνεται στη στρατηγική ελέγχου \textlatin{Q(V)} των αντιστροφέων του δικτύου, με κύριο στόχο τη ρύθμιση της τάσης καθώς  αυτή υπερβαίνει τα επιτρεπτά όρια για μεγαλά χρονικά διαστήματα. Μέσω της λειτουργίας αυτής επιτυγχάνεται δυναμική προσαρμογή της άεργου ισχύος, ενώ παράλληλα αξιολογούνται οι \textbf{επιπτώσεις} τόσο στις \textbf{απώλειες} του δικτύου όσο και στα επίπεδα \textbf{έγχυσης άεργου ισχύος}.\\
Oι υποεξέταση κόμβοι του δικτύου είναι 3: 3,47,15 οπου ειναι ο πιο κοντινός στο ζυγό και οι 2 πιο απομακρυσμένοι απο τις 2 διακλαδωσεις του δικτύου.
Παρακάτω παρατίθενται η λειτουργία τους για όλο το χρόνο .

Αρχικά γίνεται μια ανάλυση για τον κόμβο 47 γίνεται ανάλυση για όλη τη χρονιά:


\begin{figure}[H]
    \centering
    \includegraphics[width=0.90\textwidth]{PICS/term_47_Q(V)_ALL_YEAR.jpg}
    \caption{Ετήσια λειτουργία του ελέγχου \textlatin{Q(V)} στον κόμβο 47}
    \label{fig:term47_qv_all_year}
\end{figure}
Παρατηρείται πως με την εφαρμογη του ελέγχου η τάση διατηρείται μεσα στα επιτρεπτά όρια.Επιπλέον οι αλλοι 2 κόμβοι 3,15 είχαν πολυ καλύτερη συμπεριφορά ως προς τη τάση.
